\begin{abstract}
Parameter-efficient fine-tuning (PEFT) is widely used to adapt large pre-trained models, but its behavior on geospatial foundation models is easily confounded by sensor mismatch, spectral-channel bridging, domain shift, and architecture-specific implementation details. This paper studies these confounders for Prithvi-100M. We evaluate seven adaptation strategies across five heterogeneous input configurations on EuroSAT, validate the main ranking on BigEarthNet-S2, extend the protocol to LandCover.ai semantic segmentation, and test whether the ranking survives a production-style Linhe County land-use/land-cover workflow and LoveDA urban/rural cross-domain transfer. The study covers 109 public-benchmark runs, a 35{,}920-patch Linhe corpus supervised by Esri-derived land-cover labels, a 30-run LoveDA cross-domain sweep, LoveDA full fine-tuning baselines, and a completed EuroSAT channel-bridge ablation comparing deterministic zero padding with a learned 10-to-6 channel projection.

The results support a diagnostic rather than universal conclusion. First, LoRA underperformance on Prithvi-100M has two separable causes: standard module-wise insertion misses the packed query, key, and value projections in fused-QKV attention, and corrected split-QKV LoRA still saturates below bottleneck adapters even when its rank is increased to a parameter budget comparable to or larger than Houlsby adapters. Second, Houlsby adapters provide the most reliable adaptation among the evaluated PEFT methods, reaching 94.5\% of the full fine-tuning score on EuroSAT s2\_full while training 1.4\% of the parameters, and widening the advantage on dense prediction. Third, input modality and domain shift strongly condition PEFT behavior: learned channel bridging improves EuroSAT accuracy while preserving the Houlsby-over-linear ordering, and LoveDA cross-domain diagnostics suggest that sub-million-parameter adapters can recover weak binary structure without recovering full multi-class semantic transfer.

These findings do not establish a universal PEFT ranking for all geospatial foundation models. They show that architecture-aware inspection, channel-bridge controls, and parameter-capacity audits are necessary before transferring PEFT intuitions from mainstream vision to Prithvi-like remote-sensing backbones. The resulting benchmark is therefore both a practical guide for Prithvi-100M deployment and a reproducible stress-test protocol for future GeoFM adaptation studies.
\end{abstract}
